\chapter{Phân tích lỗi}
\section{Lỗi sentiment classification}
Theo \code{results/module5/error_analysis.md}, lớp neutral/mixed là lớp khó nhất. Linear SVM đạt neutral F1 khoảng 0,8230, thấp hơn negative và positive. Nguyên nhân chính là neutral ít dữ liệu hơn và nội dung neutral thường không rõ cực cảm xúc. Một review 3 sao có thể vừa khen vừa chê, ví dụ “giao nhanh nhưng chất lượng bình thường”. Nếu ánh xạ 3 sao thành neutral, mô hình có thể khó học vì từ khóa tích cực và tiêu cực cùng xuất hiện.

Một nguồn lỗi khác là rating không luôn khớp sentiment thật. Có người cho 5 sao để ủng hộ shop nhưng vẫn phàn nàn về giao hàng. Có người cho 1 sao vì sự cố cá nhân nhưng nội dung không nêu rõ vấn đề. Vì nhãn sentiment được suy ra từ rating, nhiễu nhãn là không tránh khỏi.

Review ngắn cũng gây lỗi. Câu “tạm” hoặc “ok” có thể là positive nhẹ hoặc neutral tùy ngữ cảnh. Review có mỉa mai, phủ định kép, viết tắt, sai chính tả cũng khó cho TF-IDF + SVM. Ví dụ “không phải là không tốt” có cấu trúc phủ định phức tạp mà mô hình tuyến tính khó hiểu sâu.

\section{Lỗi retrieval}
Retrieval có thể sai vì query chung chung. Query “khách hàng phàn nàn gì về sản phẩm” rất rộng, có thể truy xuất nhiều loại phàn nàn khác nhau. Nếu người dùng hỏi cụ thể hơn như “review nào nói shop giao thiếu phụ kiện”, evidence sẽ dễ chính xác hơn.

Dense retrieval có thể lấy review cùng chủ đề nhưng sai polarity. Ví dụ query hỏi phàn nàn về đóng gói nhưng kết quả lại là review khen “đóng gói kỹ”. Đây là lý do cần sentiment filter, keyword rerank hoặc hybrid retrieval. Ngược lại, BM25 có thể match keyword nhưng sai intent, ví dụ review nói “không có gì phàn nàn” nhưng vẫn chứa từ “phàn nàn”.

Review spam hoặc review nhận xu cũng làm nhiễu retrieval. Một số review dài không liên quan đến sản phẩm có thể chứa nhiều từ khóa phổ biến và được kéo vào top-k. Hệ thống có một số bước lọc và citation, nhưng chưa có mô hình lọc spam chuyên biệt.

\section{Lỗi RAG}
RAG trong đề tài là template-based. Ưu điểm là an toàn, có citation, không cần API key và ít hallucination. Hạn chế là câu trả lời kém linh hoạt hơn LLM. Nếu retrieved reviews yếu, template không thể tự suy luận đúng. Vì vậy hệ thống cần báo “Chưa đủ bằng chứng rõ ràng” khi score hoặc nội dung evidence không đủ.

Một lỗi thường gặp là lặp ý. Nếu top-k có nhiều review giống nhau, câu trả lời có thể lặp cùng một vấn đề. Cách giảm là deduplicate evidence theo comment/product hoặc gom nhóm aspect trước khi tạo answer.

\section{Lỗi Module 6}
Module 6 dùng rule-based aspect sentiment nên minh bạch nhưng chưa đủ mạnh. Nó có thể nhầm khi review chứa phủ định phức tạp, mỉa mai hoặc typo. Ví dụ “không đến nỗi tệ” có từ “tệ” nhưng sentiment không hẳn negative. Benchmark Module 6 còn nhỏ, nên kết quả \code{keyword_hit@5} bằng 1,000 không thể xem là bằng chứng tổng quát.

CrossEncoder reranker là optional. Nếu môi trường không có model hoặc không có mạng, hệ thống fallback lexical reranker. Fallback giúp chạy ổn định nhưng chất lượng có thể thấp hơn reranker neural thật.


\section{Taxonomy lỗi tổng hợp}
Bảng \ref{tab:error-taxonomy} gom các nhóm lỗi chính, nguyên nhân và cách giảm thiểu. Việc trình bày theo taxonomy giúp phần phân tích lỗi có tính hệ thống thay vì chỉ kể từng trường hợp rời rạc.
\begin{longtable}{L{0.24\textwidth}L{0.34\textwidth}L{0.32\textwidth}}
\caption{Taxonomy lỗi của hệ thống}\label{tab:error-taxonomy}\\
\toprule
Nhóm lỗi & Nguyên nhân thường gặp & Hướng giảm thiểu \\
\midrule
Label noise & Rating không khớp nội dung, 5 sao nhưng vẫn chê hoặc 1 sao vì lý do ngoài sản phẩm & Gán nhãn thủ công một tập nhỏ, dùng active learning, kiểm tra conflict duplicate \\
Neutral/mixed & Review 3 sao có cả khen và chê; số mẫu neutral ít & Thêm dữ liệu neutral, dùng model ngữ cảnh sâu, tách mixed sentiment \\
Phủ định/typo/không dấu & "không tệ", "ko dùng dc", "giao hang lau" & Chuẩn hóa text tốt hơn, bổ sung từ điển biến thể, dùng PhoBERT/XLM-R \\
Sarcasm/mỉa mai & Câu nghe tích cực nhưng ý tiêu cực & Cần dữ liệu annotated, model ngữ cảnh và human review \\
Retrieval sai polarity & Query hỏi phàn nàn nhưng lấy review khen cùng chủ đề & Sentiment filter, keyword rerank, aspect-aware rules \\
Keyword match sai intent & Có từ khóa giống query nhưng không trả lời đúng câu hỏi & Reranker supervised hoặc CrossEncoder ổn định hơn \\
Duplicate evidence & Top-k có nhiều review giống nhau hoặc cùng sản phẩm & Deduplicate theo comment/product trước khi generate answer \\
Template lặp ý & Các evidence cùng aspect làm answer lặp & Gom nhóm aspect và giới hạn citation mỗi ý \\
Module 6 rules & Rule-based aspect sentiment không xử lý tốt phủ định phức tạp & Xây dataset ABSA và huấn luyện model supervised \\
\bottomrule
\end{longtable}

\section{Bài học từ phân tích lỗi}
Phân tích lỗi cho thấy hệ thống không chỉ cần model tốt mà còn cần dữ liệu sạch, evaluation rõ ràng và giao diện không gây hiểu nhầm. Với sentiment, cần chú ý neutral và label noise. Với RAG, cần tách rõ score retrieval và relevance thật. Với giao diện thử nghiệm, cần tách sentiment input khỏi query RAG. Những bài học này đã được phản ánh trong phiên bản giao diện cuối cùng.
